\paragraph{Weyl 不变量像、双框架矩与黄金分支证书振荡}
在 window-\(6\) 的 \(B_3/C_3\) 根字典、\(C_3\) 二次矩各向同性、\(B_3/C_3\) 线性统一障碍以及黄金分支的 Kronecker--Wasserstein 奇偶证书已经建立之后，还可进一步把 Weyl 不变量、方向矩、有限框架与 Fibonacci 归一化误差证书压缩为同一组闭合后果。

\begin{theorem}[window-\texorpdfstring{$6$}{6} 循环根字典上的 Weyl 不变量像恰由平方范数生成]\label{thm:xi-window6-b3c3-weyl-invariant-image-quadratic}
沿用定理 \ref{thm:xi-window6-b3c3-no-global-linear-unifier} 的记号，并记
\[
R_B:=\iota_B(X_6^{\mathrm{cyc}}),
\qquad
R_C:=\iota_C(X_6^{\mathrm{cyc}}).
\]
分别写
\[
R_{B,\mathrm{sh}}:=\{\pm e_1,\pm e_2,\pm e_3\},
\qquad
R_{B,\mathrm{lg}}:=\{\pm e_i\pm e_j:\ 1\le i<j\le 3\},
\]
\[
R_{C,\mathrm{sh}}:=\{\pm e_i\pm e_j:\ 1\le i<j\le 3\},
\qquad
R_{C,\mathrm{lg}}:=\{\pm 2e_1,\pm 2e_2,\pm 2e_3\}.
\]
令
\[
W_B:=W(B_3)\cong (\ZZ/2\ZZ)^3\rtimes S_3,
\qquad
W_C:=W(C_3)\cong (\ZZ/2\ZZ)^3\rtimes S_3.
\]
则：
\begin{enumerate}
  \item \(W_B\) 在 \(R_B\) 上恰有两条轨道 \(R_{B,\mathrm{sh}}\) 与 \(R_{B,\mathrm{lg}}\)，其大小分别为 \(6\) 与 \(12\)；\(W_C\) 在 \(R_C\) 上同样恰有两条轨道 \(R_{C,\mathrm{sh}}\) 与 \(R_{C,\mathrm{lg}}\)，其大小亦分别为 \(12\) 与 \(6\)；
  \item 限制映射
  \[
  \operatorname{res}_B:\mathbb R[x_1,x_2,x_3]^{W_B}\to \mathbb R^{R_B},
  \qquad
  \operatorname{res}_C:\mathbb R[x_1,x_2,x_3]^{W_C}\to \mathbb R^{R_C}
  \]
  的像都恰为二维，并且都由常数函数 \(1\) 与平方范数
  \[
  Q(x):=|x|^2
  \]
  的限制生成；
  \item 在 \(B_3\) 规范下，短根与长根的指示函数满足
  \[
  \boxed{\;
  \mathbf 1_{R_{B,\mathrm{sh}}}(x)=2-|x|^2,
  \qquad
  \mathbf 1_{R_{B,\mathrm{lg}}}(x)=|x|^2-1
  \qquad(x\in R_B).\;}
  \]
  \item 在 \(C_3\) 规范下，相应地有
  \[
  \boxed{\;
  \mathbf 1_{R_{C,\mathrm{sh}}}(x)=\frac{4-|x|^2}{2},
  \qquad
  \mathbf 1_{R_{C,\mathrm{lg}}}(x)=\frac{|x|^2-2}{2}
  \qquad(x\in R_C).\;}
  \]
\end{enumerate}
\end{theorem}
\begin{proof}
由表 \ref{tab:fold6_b3c3_root_dictionary_b3}，\(R_B\) 正是
\[
\{\pm e_1,\pm e_2,\pm e_3\}\cup \{\pm e_i\pm e_j:\ 1\le i<j\le 3\},
\]
而由表 \ref{tab:fold6_b3c3_root_dictionary_c3}，\(R_C\) 正是
\[
\{\pm 2e_1,\pm 2e_2,\pm 2e_3\}\cup \{\pm e_i\pm e_j:\ 1\le i<j\le 3\}.
\]
群 \(W_B\) 与 \(W_C\) 都由坐标置换与独立符号翻转生成，因此分别在各自的轴向根集合与非轴向根集合上作用传递，于是两侧都恰有两条轨道，大小由集合基数直接读出。

任一 Weyl 不变多项式在有限根集上的限制都必须在每条 Weyl 轨道上取常值，因此 \(\operatorname{im}(\operatorname{res}_B)\) 与 \(\operatorname{im}(\operatorname{res}_C)\) 的维数都至多为 \(2\)。另一方面，在 \(R_B\) 上
\[
Q(x)=|x|^2=
\begin{cases}
1,& x\in R_{B,\mathrm{sh}},\\
2,& x\in R_{B,\mathrm{lg}},
\end{cases}
\]
而在 \(R_C\) 上
\[
Q(x)=|x|^2=
\begin{cases}
2,& x\in R_{C,\mathrm{sh}},\\
4,& x\in R_{C,\mathrm{lg}}.
\end{cases}
\]
故 \(1\) 与 \(Q\) 在两条轨道上分别取不同值，从而在线性上独立，并生成全部轨道常值函数空间。于是两侧像空间都恰为二维。

最后，只需对两条轨道上的取值解线性插值方程即可。在 \(R_B\) 上，
\[
a+b|x|^2=
\begin{cases}
1,& |x|^2=1,\\
0,& |x|^2=2
\end{cases}
\]
解得 \(a=2,b=-1\)，从而
\[
\mathbf 1_{R_{B,\mathrm{sh}}}(x)=2-|x|^2.
\]
互补轨道的指示函数即为
\[
\mathbf 1_{R_{B,\mathrm{lg}}}(x)=1-\mathbf 1_{R_{B,\mathrm{sh}}}(x)=|x|^2-1.
\]
在 \(R_C\) 上同理，对
\[
a+b|x|^2=
\begin{cases}
1,& |x|^2=2,\\
0,& |x|^2=4
\end{cases}
\]
求解得到 \(a=2,b=-\frac12\)，故
\[
\mathbf 1_{R_{C,\mathrm{sh}}}(x)=\frac{4-|x|^2}{2},
\qquad
\mathbf 1_{R_{C,\mathrm{lg}}}(x)=1-\mathbf 1_{R_{C,\mathrm{sh}}}(x)=\frac{|x|^2-2}{2}.
\]
证毕。
\end{proof}
\leanverified{paper\_xi\_window6\_b3c3\_weyl\_invariant\_image\_quadratic}

\begin{theorem}[window-\texorpdfstring{$6$}{6} 的 \texorpdfstring{$B_3/C_3$}{B3/C3} 字典形成双紧框架且不存在单一 Euclid 相似统一]\label{thm:xi-window6-b3c3-tight-frame-fourth-moment-nonsimilarity}
沿用定理 \ref{thm:xi-window6-b3c3-weyl-invariant-image-quadratic} 的记号。则二阶矩算子满足
\[
\boxed{\;
\sum_{r\in R_B}rr^{\mathsf T}=10I_3,
\qquad
\sum_{r\in R_C}rr^{\mathsf T}=16I_3.\;}
\]
因而 \(R_B\) 与 \(R_C\) 分别构成 \(\mathbb R^3\) 上框架常数为 \(10\) 与 \(16\) 的有限紧框架。

进一步，对任意 \(v=(v_1,v_2,v_3)\in\mathbb R^3\)，四阶方向矩满足
\[
\boxed{\;
\sum_{r\in R_B}\langle r,v\rangle^4
=
10\sum_{i=1}^3 v_i^4
+24\sum_{1\le i<j\le 3}v_i^2v_j^2,\;}
\]
\[
\boxed{\;
\sum_{r\in R_C}\langle r,v\rangle^4
=
40\sum_{i=1}^3 v_i^4
+24\sum_{1\le i<j\le 3}v_i^2v_j^2.\;}
\]
最后，不存在 \(s>0\) 与 \(O\in O(3)\) 使
\[
R_C=s\,O(R_B).
\]
换言之，\(B_3/C_3\) 两套字典之间的长短根互换并非单一 Euclid 相似变换，而必然是按 Weyl 轨道分段的重标定。
\leanverified{paper\_xi\_window6\_b3c3\_tight\_frame\_fourth\_moment\_nonsimilarity}
\end{theorem}
\begin{proof}
对 \(R_B\)，六个短根 \(\pm e_i\) 的贡献为
\[
\sum_{i=1}^3\bigl(e_ie_i^{\mathsf T}+(-e_i)(-e_i)^{\mathsf T}\bigr)=2I_3.
\]
而对每一对指标 \(1\le i<j\le 3\)，四个长根 \(\pm e_i\pm e_j\) 的和满足
\[
\sum_{\varepsilon_i,\varepsilon_j\in\{\pm1\}}
(\varepsilon_i e_i+\varepsilon_j e_j)(\varepsilon_i e_i+\varepsilon_j e_j)^{\mathsf T}
=
4(e_ie_i^{\mathsf T}+e_je_j^{\mathsf T}),
\]
因为交叉项按符号对称抵消。对三对 \((i,j)\) 求和便得到总贡献 \(8I_3\)，从而
\[
\sum_{r\in R_B}rr^{\mathsf T}=2I_3+8I_3=10I_3.
\]

对 \(R_C\)，六个长根 \(\pm2e_i\) 的贡献为
\[
\sum_{i=1}^3\bigl((2e_i)(2e_i)^{\mathsf T}+(-2e_i)(-2e_i)^{\mathsf T}\bigr)=8I_3,
\]
而十二个短根 \(\pm e_i\pm e_j\) 的贡献与上面完全相同，亦为 \(8I_3\)。故
\[
\sum_{r\in R_C}rr^{\mathsf T}=16I_3.
\]
于是对任意 \(v\in\mathbb R^3\)，都有
\[
\sum_{r\in R_B}\langle r,v\rangle^2=10\|v\|_2^2,
\qquad
\sum_{r\in R_C}\langle r,v\rangle^2=16\|v\|_2^2,
\]
即二者分别为框架常数 \(10\) 与 \(16\) 的有限紧框架。

再看四阶矩。对 \(R_B\) 的六个短根，
\[
\sum_{r\in R_{B,\mathrm{sh}}}\langle r,v\rangle^4=2\sum_{i=1}^3 v_i^4.
\]
而对任一固定对 \((i,j)\)，有恒等式
\[
\sum_{\varepsilon_i,\varepsilon_j\in\{\pm1\}}
(\varepsilon_i v_i+\varepsilon_j v_j)^4
=
4(v_i^4+6v_i^2v_j^2+v_j^4).
\]
对三对 \((i,j)\) 求和便得到长根部分贡献
\[
8\sum_{i=1}^3 v_i^4+24\sum_{1\le i<j\le 3}v_i^2v_j^2.
\]
两者相加即得
\[
\sum_{r\in R_B}\langle r,v\rangle^4
=
10\sum_{i=1}^3 v_i^4+24\sum_{1\le i<j\le 3}v_i^2v_j^2.
\]
对 \(R_C\)，长根 \(\pm 2e_i\) 部分给出
\[
\sum_{r\in R_{C,\mathrm{lg}}}\langle r,v\rangle^4=32\sum_{i=1}^3 v_i^4,
\]
而十二个短根仍给出
\[
8\sum_{i=1}^3 v_i^4+24\sum_{1\le i<j\le 3}v_i^2v_j^2,
\]
故
\[
\sum_{r\in R_C}\langle r,v\rangle^4
=
40\sum_{i=1}^3 v_i^4+24\sum_{1\le i<j\le 3}v_i^2v_j^2.
\]

最后，若存在 \(s>0\) 与 \(O\in O(3)\) 使 \(R_C=s\,O(R_B)\)，则由二阶矩公式必有
\[
16I_3
=
\sum_{r\in R_C}rr^{\mathsf T}
=
s^2 O\left(\sum_{r\in R_B}rr^{\mathsf T}\right)O^{\mathsf T}
=
10s^2I_3,
\]
故
\[
s^2=\frac85,
\qquad
s^4=\frac{64}{25}.
\]
对任意单位向量 \(v\)，记 \(x_i:=v_i^2\)，则 \(\sum_i x_i=1\)。上面的四阶矩公式给出
\[
\sum_{r\in R_B}\langle r,v\rangle^4
=
10\sum_i x_i^2+24\sum_{i<j}x_ix_j
=
12-2\sum_i x_i^2,
\]
故其取值范围为
\[
\left[10,\frac{34}{3}\right],
\]
因为 \(\sum_i x_i^2\in[1/3,1]\)。若乘上 \(s^4=64/25\)，则可能范围变为
\[
\left[\frac{128}{5},\frac{2176}{75}\right].
\]
而对 \(R_C\) 同理，
\[
\sum_{r\in R_C}\langle r,v\rangle^4
=
40\sum_i x_i^2+24\sum_{i<j}x_ix_j
=
12+28\sum_i x_i^2,
\]
其取值范围为
\[
\left[\frac{64}{3},40\right].
\]
两者范围不同，因此不可能满足
\[
\sum_{r\in R_C}\langle r,v\rangle^4
=
s^4\sum_{r\in R_B}\langle r,O^{\mathsf T}v\rangle^4
\qquad(\forall\,\|v\|_2=1).
\]
故这样的 Euclid 相似变换不存在。证毕。
\end{proof}

\begin{corollary}[黄金分支在 Fibonacci 分母上的星差与传输证书具有严格奇偶分裂]\label{cor:xi-golden-fibonacci-star-w1-normalized-parity}
沿用定理 \ref{thm:xi-golden-w1-star-certificate-parity-splitting} 的记号。记
\[
D_n^\ast:=D_{F_n}^\ast(P_{F_n}(\varphi^{-1})),
\qquad
\mathsf{T}_n:=\mathsf{T}_{F_n}(\varphi^{-1}).
\]
则 Fibonacci 归一化的星差满足
\[
\boxed{\;
F_nD_n^\ast=
\begin{cases}
1+\dfrac{1}{\sqrt5}+o(1),& n\ \mathrm{even},\\[2mm]
1,& n\ \mathrm{odd},
\end{cases}}
\]
其中奇数支是精确恒等式。

进一步，
\[
\boxed{\;
\frac{\mathsf{T}_n}{D_n^\ast}\to
\begin{cases}
\dfrac12,& n\ \mathrm{even},\\[2mm]
\dfrac12-\dfrac{1}{2\sqrt5}+\dfrac1{15},& n\ \mathrm{odd}.
\end{cases}}
\]
特别地，好侧奇数子序列上保留下来的额外箱内几何常数正是
\[
\frac12-\frac{1}{2\sqrt5}+\frac1{15}.
\]
\end{corollary}
\begin{proof}
由定理 \ref{thm:xi-golden-w1-star-certificate-parity-splitting}，当 \(n\) 为偶数时有精确恒等式
\[
\mathsf{T}_n=\frac12 D_n^\ast.
\]
再由定理 \ref{thm:xi-golden-w1-true-two-phase-limit} 的偶数支，
\[
F_n\mathsf{T}_n\to \frac12+\frac{1}{2\sqrt5},
\]
故
\[
F_nD_n^\ast=2F_n\mathsf{T}_n\to 1+\frac1{\sqrt5}.
\]

当 \(n\) 为奇数时，由结论 \ref{con:xi-terminal-kronecker-w1-denominator-one-sided-sensitivity}，黄金分支处于好侧，因而
\[
D_n^\ast=\frac1{F_n},
\]
从而 \(F_nD_n^\ast=1\) 为精确恒等式。再由定理
\ref{thm:xi-golden-w1-star-certificate-parity-splitting} 的奇数支，
\[
\frac{\mathsf{T}_n}{D_n^\ast}
=
F_n\mathsf{T}_n
\to
\frac12-\frac{1}{2\sqrt5}+\frac1{15}.
\]
证毕。
\end{proof}

\begin{corollary}[黄金分支 Fibonacci 归一化证书的振幅与 Ces\`aro 常数]\label{cor:xi-golden-fibonacci-star-w1-oscillation-cesaro}
\leanverified{paper\_xi\_golden\_fibonacci\_star\_w1\_oscillation\_cesaro}
沿用推论 \ref{cor:xi-golden-fibonacci-star-w1-normalized-parity} 的记号。则 Fibonacci 归一化传输差异与星差的振幅分别为
\[
\boxed{\;
\limsup_{n\to\infty}F_n\mathsf{T}_n
-
\liminf_{n\to\infty}F_n\mathsf{T}_n
=
\frac1{\sqrt5}-\frac1{15},\;}
\]
\[
\boxed{\;
\limsup_{n\to\infty}F_nD_n^\ast
-
\liminf_{n\to\infty}F_nD_n^\ast
=
\frac1{\sqrt5}.\;}
\]
并且 Ces\`aro 平均满足
\[
\boxed{\;
\lim_{N\to\infty}\frac1N\sum_{n=1}^N F_n\mathsf{T}_n
=
\frac{8}{15},\;}
\]
\[
\boxed{\;
\lim_{N\to\infty}\frac1N\sum_{n=1}^N F_nD_n^\ast
=
1+\frac{1}{2\sqrt5}.\;}
\]
\end{corollary}
\begin{proof}
传输差异的振幅公式与 Ces\`aro 平均公式分别就是推论
\ref{cor:xi-golden-w1-two-phase-amplitude-center} 与
\ref{cor:xi-golden-w1-two-phase-arithmetic-mean-eight-fifteenths} 的内容。

对星差，由推论 \ref{cor:xi-golden-fibonacci-star-w1-normalized-parity}，
\[
\limsup_{n\to\infty}F_nD_n^\ast=1+\frac1{\sqrt5},
\qquad
\liminf_{n\to\infty}F_nD_n^\ast=1,
\]
相减即得振幅
\[
\frac1{\sqrt5}.
\]

同样，由偶数与奇数索引在自然数中各占密度 \(1/2\)，且对应两条有界子序列分别收敛到
\[
1+\frac1{\sqrt5},
\qquad
1,
\]
故其 Ces\`aro 平均收敛到二者的算术平均
\[
\frac12\left(1+\frac1{\sqrt5}\right)+\frac12\cdot 1
=
1+\frac{1}{2\sqrt5}.
\]
证毕。
\end{proof}

\begin{theorem}[window-\texorpdfstring{$6$}{6} 循环根字典上的 Weyl 平均概率测度完全由单一二次矩恢复]\label{thm:xi-window6-b3c3-weyl-average-measure-second-moment}
沿用定理 \ref{thm:xi-window6-b3c3-weyl-invariant-image-quadratic} 的记号。令
\[
\sigma_{B,\mathrm{sh}},\ \sigma_{B,\mathrm{lg}}
\]
分别为 \(R_{B,\mathrm{sh}}\)、\(R_{B,\mathrm{lg}}\) 上的均匀概率测度，令
\[
\sigma_{C,\mathrm{sh}},\ \sigma_{C,\mathrm{lg}}
\]
分别为 \(R_{C,\mathrm{sh}}\)、\(R_{C,\mathrm{lg}}\) 上的均匀概率测度。对任意支撑在 \(R_B\) 上的概率测度 \(\nu_B\)，定义其 Weyl 平均为
\[
\overline{\nu}_B:=\frac{1}{|W_B|}\sum_{g\in W_B}g_\ast\nu_B,
\]
并记
\[
m_2^{(B)}(\nu_B):=\int_{R_B}|x|^2\,d\nu_B(x).
\]
则有显式恢复公式
\[
\boxed{\;
\overline{\nu}_B
=
\bigl(2-m_2^{(B)}(\nu_B)\bigr)\sigma_{B,\mathrm{sh}}
+
\bigl(m_2^{(B)}(\nu_B)-1\bigr)\sigma_{B,\mathrm{lg}}.\;}
\]

同样，对任意支撑在 \(R_C\) 上的概率测度 \(\nu_C\)，定义
\[
\overline{\nu}_C:=\frac{1}{|W_C|}\sum_{g\in W_C}g_\ast\nu_C,
\qquad
m_2^{(C)}(\nu_C):=\int_{R_C}|x|^2\,d\nu_C(x),
\]
则有
\[
\boxed{\;
\overline{\nu}_C
=
\frac{4-m_2^{(C)}(\nu_C)}{2}\,\sigma_{C,\mathrm{sh}}
+
\frac{m_2^{(C)}(\nu_C)-2}{2}\,\sigma_{C,\mathrm{lg}}.\;}
\]
因此，\(R_B\) 与 \(R_C\) 上全部 Weyl 不变概率测度构成的单纯形都退化为一条闭线段，其唯一自由参数就是二次矩。
\leanverified{paper\_xi\_window6\_b3c3\_weyl\_average\_measure\_second\_moment}
\end{theorem}
\begin{proof}
对 \(R_B\)，任意 Weyl 平均测度 \(\overline{\nu}_B\) 都在两条 Weyl 轨道上取常值，故必可唯一写成
\[
\overline{\nu}_B=a\,\sigma_{B,\mathrm{sh}}+(1-a)\,\sigma_{B,\mathrm{lg}}
\qquad (0\le a\le 1).
\]
由于 \(|x|^2\) 是 \(W_B\)-不变函数，故
\[
m_2^{(B)}(\nu_B)=m_2^{(B)}(\overline{\nu}_B).
\]
再利用 \(R_{B,\mathrm{sh}}\) 上 \(|x|^2=1\)、\(R_{B,\mathrm{lg}}\) 上 \(|x|^2=2\)，得到
\[
m_2^{(B)}(\nu_B)
=
a\cdot 1+(1-a)\cdot 2
=
2-a.
\]
解得
\[
a=2-m_2^{(B)}(\nu_B),
\qquad
1-a=m_2^{(B)}(\nu_B)-1,
\]
从而得到第一条公式。

对 \(R_C\) 同理，任意 Weyl 平均都唯一写成
\[
\overline{\nu}_C=b\,\sigma_{C,\mathrm{sh}}+(1-b)\,\sigma_{C,\mathrm{lg}},
\]
并由 \(R_{C,\mathrm{sh}}\) 上 \(|x|^2=2\)、\(R_{C,\mathrm{lg}}\) 上 \(|x|^2=4\) 得
\[
m_2^{(C)}(\nu_C)=m_2^{(C)}(\overline{\nu}_C)=2b+4(1-b)=4-2b.
\]
故
\[
b=\frac{4-m_2^{(C)}(\nu_C)}{2},
\qquad
1-b=\frac{m_2^{(C)}(\nu_C)-2}{2},
\]
即得第二条公式。

最后，Weyl 不变概率测度恰是两轨道均匀测度的凸组合，因此其参数空间正是一条闭线段，而上式表明二次矩已经给出这一线段上的唯一坐标。证毕。
\end{proof}

\noindent
由此，window-\(6\) 的 \(B_3/C_3\) 根字典不再只是一个离散配对表，而在函数层、矩层与概率层都表现出同一条低维刚性：Weyl 不变量的限制像由单一二次函数生成，二组根字典分别形成显式可审计的紧框架而彼此又不能由单一 Euclid 相似统一，黄金分支的 Fibonacci 证书则在星差与传输之间呈现严格奇偶振荡，而对称概率信息最终又坍缩到单一二次矩这一唯一自由度。

\endinput
